\documentclass[11pt]{article}
\usepackage[margin=1in]{geometry}
\usepackage{amsmath,amssymb,amsthm}
\usepackage{graphicx}
\usepackage[hidelinks]{hyperref}

\newtheorem{theorem}{Theorem}
\newtheorem{corollary}[theorem]{Corollary}
\newtheorem{lemma}[theorem]{Lemma}
\newcommand{\R}{\mathbb R}
\newcommand{\E}{\mathbb E}
\newcommand{\Q}{\mathcal Q}

\emergencystretch=2em

\begin{document}
\pagestyle{plain}

\begin{center}
{\Large Full local stability of the Gaussian propeller}\\[0.35em]
{\small Run \texttt{fa\_banach\_001}, arXiv:1112.2993, partial result}
\end{center}

\section*{Source problem}

For a centered finite list $v_1,\ldots,v_k$ in a real Hilbert space, let
$G$ be standard Gaussian on its span and set
\[
 \Q(v_1,\ldots,v_k)
 =\frac{\bigl(\E\max_{1\le i\le k}\langle G,v_i\rangle\bigr)^2}
 {\sum_{i=1}^k\|v_i\|^2}.
\]
The Gaussian formulation of the Propeller Conjecture asks whether
\begin{equation}\label{eq:global}
 \Q(v_1,\ldots,v_k)\le c_0,
 \qquad c_0=\frac{9}{8\pi},
\end{equation}
for every $k$ and every ambient dimension; see Khot--Naor
\cite{KN} and Heilman--Jagannath--Naor \cite{HJN}.  The latter paper proves
the three-dimensional case. On page 7 it explicitly identifies a proof
extending to all dimensions as the next natural step.

\begin{center}
\includegraphics[width=0.98\linewidth]{figures/open_problem_crop.png}
\end{center}

The global problem remains open. We prove that the conjectured optimizer is
locally rigid under arbitrary higher-dimensional perturbations, substantially
strengthening the previously isolated purely transverse case.

\section*{Proof intuition}

Choose the three vertices that are close to the propeller blades and translate
the entire list by their mean. Translation does not change the expected
Gaussian maximum. The translated list consists of a centered active triangle
and a small arbitrary residual cloud. Re-centering the whole list gives an
exact residual denominator cost of order $\eta^2$.

The support function of a nondegenerate centered triangle dominates a fixed
multiple of the planar radius. A residual vertex of total size $\eta$ can win
only when the planar Gaussian component is unusually small compared with its
orthogonal component. The small-ball probability of a two-dimensional
Gaussian is quadratic, so the expected improvement is only $O(\eta^3)$.
Thus the quadratic denominator cost beats the cubic numerator gain.

The active triangle needs no external sharp inequality. Cauchy's planar
perimeter formula expresses its expected Gaussian support in terms of its
perimeter, and one application of Cauchy--Schwarz gives the sharp constant
$9/(8\pi)$, with equality exactly for an equilateral triangle.

\section*{Quantitative theorem}

The following statement is useful beyond a neighborhood of the equilateral
triangle. It gives an explicit certificate for adjoining an arbitrary small
cloud to any nondegenerate centered triangle.

\begin{theorem}[Arbitrary-cloud certificate]\label{thm:certificate}
Let $u_1,u_2,u_3$ lie in a finite-dimensional real Hilbert space, satisfy
$u_1+u_2+u_3=0$, and span a two-dimensional space $U$. Put
\[
 S=\sum_{i=1}^3\|u_i\|^2,
 \quad
 M=\max_{1\le i\le3}\langle X,u_i\rangle,
 \quad
 W=\E M,
\]
where $X$ is standard Gaussian on $U$, and define the inradius parameter
\[
 \alpha=\min_{x\in U,\ \|x\|=1}\max_{1\le i\le3}\langle x,u_i\rangle>0.
\]
Fix $m\ge1$, set $k=m+3$, and let $y_1,\ldots,y_m$ be arbitrary vectors in
the ambient Hilbert space. Write
\[
 \eta^2=\sum_{j=1}^m\|y_j\|^2,
 \qquad
 A=\eta^2-\frac1k\left\|\sum_{j=1}^m y_j\right\|^2,
 \qquad
 \bar y=\frac1k\sum_{j=1}^m y_j.
\]
Form the centered list
\[
 v_i=u_i-\bar y\quad(1\le i\le3),
 \qquad
 v_{3+j}=y_j-\bar y\quad(1\le j\le m).
\]
If $0<\eta<\alpha$, define
\[
 b=\frac{\sqrt{2/\pi}}{2(\alpha-\eta)^2}.
\]
Whenever
\begin{equation}\label{eq:certificate}
 2Wb\eta^3+b^2\eta^6<c_0A,
\end{equation}
one has
\[
 \Q(v_1,\ldots,v_k)<c_0.
\]
Moreover $A\ge(3/k)\eta^2$, so \eqref{eq:certificate} holds for all
sufficiently small nonzero $\eta$, uniformly over the directions and number
of dimensions used by the residual cloud.
\end{theorem}

\begin{corollary}[Full local stability]\label{cor:local}
Fix $k\ge4$ and an equilateral centered triple
$u_1^0,u_2^0,u_3^0$. In the space of centered $k$-tuples there is a
neighborhood of
\[
 (u_1^0,u_2^0,u_3^0,0,\ldots,0)
\]
such that every tuple in this neighborhood satisfies \eqref{eq:global}.
Equality holds precisely when, after the same labeling, the first three
vectors form an equilateral centered triple and every remaining vector is
zero. Thus, modulo rotations, scaling, and relabeling, the propeller is a
strict local maximizer against all perturbations in arbitrary ambient
dimension.
\end{corollary}

\section*{Proof}

We first record the sharp estimate for the active triangle.

\begin{lemma}[Centered-triangle inequality]\label{lem:triangle}
With the notation of Theorem~\ref{thm:certificate},
\[
 W^2\le c_0S,
\]
with equality if and only if the triangle with vertices $u_1,u_2,u_3$ is
equilateral.
\end{lemma}

\begin{proof}
Let $P$ be the perimeter of $K=\operatorname{conv}\{u_1,u_2,u_3\}$.
Write $X=R\Theta$, where $\Theta$ is uniform on the unit circle in $U$ and
$R$ is independent of $\Theta$. The planar Cauchy perimeter formula says
that the angular average of the support function $h_K$ is $P/(2\pi)$.
Since $\E R=\sqrt{\pi/2}$,
\[
 W=\E h_K(X)=\frac{P}{2\sqrt{2\pi}}.
\]
If $a,b,c$ are the three side lengths, the centroid identity gives
\[
 a^2+b^2+c^2=3S.
\]
Consequently
\[
 P^2=(a+b+c)^2\le3(a^2+b^2+c^2)=9S,
\]
and therefore $W^2\le9S/(8\pi)=c_0S$. Equality in Cauchy--Schwarz occurs
exactly when $a=b=c$.
\end{proof}

\begin{proof}[Proof of Theorem~\ref{thm:certificate}]
The list $v_1,\ldots,v_k$ is centered. Adding the common centered Gaussian
$-\langle G,\bar y\rangle$ does not change the expectation of the maximum,
so
\begin{equation}\label{eq:numeratorexact}
 \E\max_{1\le i\le k}\langle G,v_i\rangle
 =\E\max\left\{M,\max_{1\le j\le m}\langle G,y_j\rangle\right\}.
\end{equation}
The denominator has the exact identity
\begin{equation}\label{eq:denomexact}
 \sum_{i=1}^k\|v_i\|^2=S+A.
\end{equation}
Indeed, this is the standard variance identity after subtracting the mean
$\bar y$. Since $\|\sum_jy_j\|^2\le m\eta^2$ and $k=m+3$,
\begin{equation}\label{eq:Alower}
 A\ge\left(1-\frac{m}{k}\right)\eta^2=\frac3k\eta^2.
\end{equation}

Decompose $G=X+Z$ according to $U\oplus U^\perp$, and write
$y_j=p_j+q_j$ in the same decomposition. Set
\[
 Q(Z)=\max_{1\le j\le m}\langle Z,q_j\rangle.
\]
Because $\|p_j\|\le\eta$ and $M\ge\alpha\|X\|$, the improvement over the
active triangle is bounded by
\begin{align}
 \delta
 &:=\E\left(\max\left\{M,\max_j\langle G,y_j\rangle\right\}-M\right)
 \notag\\
 &\le\E\bigl(Q(Z)-(\alpha-\eta)\|X\|\bigr)_+.
 \label{eq:gainreduce}
\end{align}
For $q>0$ and $a=\alpha-\eta>0$, the radius of a standard
two-dimensional Gaussian satisfies
\[
 \Pr\{\|X\|\le t\}=1-e^{-t^2/2}\le\frac{t^2}{2}.
\]
Conditioning on $Z$ therefore gives
\[
 \E_X(q-a\|X\|)_+
 \le q\Pr\{\|X\|\le q/a\}
 \le\frac{q^3}{2a^2}.
\]
Also
\[
 \E Q(Z)_+^3
 \le\sum_{j=1}^m\E\langle Z,q_j\rangle_+^3
 =\sqrt{\frac2\pi}\sum_{j=1}^m\|q_j\|^3
 \le\sqrt{\frac2\pi}\,\eta^3.
\]
Combining this with \eqref{eq:gainreduce} proves the uniform cubic bound
\begin{equation}\label{eq:cubic}
 \delta\le b\eta^3.
\end{equation}

By Lemma~\ref{lem:triangle}, \eqref{eq:numeratorexact},
\eqref{eq:denomexact}, and \eqref{eq:cubic},
\[
 \Q(v_1,\ldots,v_k)
 \le\frac{(W+b\eta^3)^2}{S+A}.
\]
The numerator on the right is
$W^2+2Wb\eta^3+b^2\eta^6$.  Lemma~\ref{lem:triangle} bounds the first
term by $c_0S$, while \eqref{eq:certificate} bounds the last two strictly
by $c_0A$. This proves the claim. The final assertion follows from
\eqref{eq:Alower}, since the left side of \eqref{eq:certificate} is
$O(\eta^3)$ and its right side is bounded below by a positive multiple of
$\eta^2$.
\end{proof}

\begin{proof}[Proof of Corollary~\ref{cor:local}]
For a centered $k$-tuple $(z_1,\ldots,z_k)$ near the displayed propeller,
let $a=(z_1+z_2+z_3)/3$, put $u_i=z_i-a$ for $i\le3$, and put
$y_j=z_{3+j}-a$. Subtracting $a$ from the entire list does not change the
expected maximum. The centered version of the resulting list is exactly the
original tuple, so Theorem~\ref{thm:certificate} applies.

In a sufficiently small neighborhood of the fixed equilateral triple, the
parameters $\alpha$ are bounded below by some $\alpha_*>0$ and the values
$W$ are bounded above by some $W_*<\infty$. If $\eta\le\alpha_*/2$, then
\[
 b\le B_*:=\frac{2\sqrt{2/\pi}}{\alpha_*^2}.
\]
Hence \eqref{eq:certificate} follows uniformly once
\[
 2W_*B_*\eta+B_*^2\eta^4<\frac{3c_0}{k},
\]
which holds for all sufficiently small positive $\eta$. If $\eta>0$, the
conclusion is strict. If $\eta=0$, all residual vectors vanish after the
translation, and centeredness forces $a=0$; Lemma~\ref{lem:triangle} gives
equality exactly for an equilateral active triple. This proves the stated
local rigidity and equality characterization.
\end{proof}

\section*{Verification notes}

The proof is analytic and has no computational dependency. The two points
most worth checking are the exact centered-denominator identity
\eqref{eq:denomexact} and the passage from \eqref{eq:gainreduce} to the
cubic estimate \eqref{eq:cubic}. Both constants are deliberately kept
explicit. The argument is dimension-free after the planar small-ball step:
the estimate
$\sum_j\|q_j\|^3\le(\sum_j\|q_j\|^2)^{3/2}$ removes any dependence on the
number of transverse directions.

As a sanity check, for an equilateral triangle of circumradius $r$ one has
$\alpha=r/2$, $S=3r^2$, and
$W=3\sqrt3\,r/(2\sqrt{2\pi})$, which gives $W^2/S=9/(8\pi)$.

\section*{Scope and novelty check}

This is a partial result, not a proof of the global higher-dimensional
conjecture. It controls every sufficiently small perturbation of the
three-blade optimizer for each fixed $k$, but it does not control
configurations far from that optimizer.

A bounded primary-source search on 2026-08-09 used the exact phrases
``Propeller Conjecture'', ``higher dimensions'', ``Gaussian moments'', and
``local stability'', together with arXiv:0807.4626 and arXiv:1112.2993.
It also inspected the primary arXiv records 1901.03934 and 2209.11216. The
search found the known three-dimensional theorem and later conditional work,
but no unconditional all-dimensional solution and no theorem matching the
arbitrary-perturbation local statement above. Novelty confidence is moderate.

The earlier run packet
\texttt{0807.4626\_propeller\_transverse\_local\_stability} proves only the
special case where a fixed residual cloud is scaled in a space orthogonal to
an undeformed equilateral triangle. The present theorem removes both
restrictions and supplies a quantitative certificate. Human review is
recommended, with priority on \eqref{eq:denomexact} and
\eqref{eq:cubic}.

\begin{thebibliography}{9}
\bibitem{KN}
S. Khot and A. Naor,
\emph{Approximate kernel clustering},
arXiv:0807.4626, 2008.

\bibitem{HJN}
S. Heilman, A. Jagannath, and A. Naor,
\emph{Solution of the propeller conjecture in $\R^3$},
arXiv:1112.2993, 2011.
\end{thebibliography}

\end{document}
